% Source: https://artofproblemsolving.com/community/c3h2182438_polynomials__binomial_theorem
% Stable identity: aops-topic-2182438-q01
The slope of $f$ is $\dfrac{18-6}{8-2}=2$, so $f(x)=6+2(x-2)=2x+2$. Then
\[q(x)=(2x+2)^{10}=2^{10}(x+1)^{10}.\]

\textbf{(a)} $q(x)=0$ exactly when $2x+2=0$, so the only solution (real or complex) is $x=-1$, and the sum of the distinct solutions is $\boxed{-1}$.

\textbf{(b)} The root $-1$ has multiplicity $10$, so the sum of the roots counted with multiplicity is $10\cdot(-1)=\boxed{-10}$. This agrees with Vieta's formulas: by the binomial theorem $q(x)=1024x^{10}+10\cdot1024x^9+\cdots$, and the sum of the roots is $-\dfrac{10240}{1024}=-10$.
